% !TeX root = ../navier-stokes-manuscript.tex
% ============================================================
% 2.1 Vortex Stretching
% ============================================================

The momentum equation has four pieces:
\[
  \underbrace{\partial_t u}_{\text{local acceleration}}
  +
  \underbrace{(u\cdot\nabla)u}_{\text{nonlinear transport}}
  =
  \underbrace{-\nabla p}_{\text{pressure}}
  +
  \underbrace{\nu\Delta u}_{\text{viscous diffusion}}.
\]
The central regularity problem lies in the competition between nonlinear transport and viscosity.
The hope is simple: viscosity smooths the fluid faster than nonlinear transport can sharpen it.
If that remains true at every scale, the velocity stays smooth.

\subsection{Vortex Stretching}\label{sub:ThePerfectlyStretchingVortex}

Take a small material fluid parcel rotating around an axis. In the local axisymmetric aligned-strain model, axial stretching lengthens the parcel while incompressibility contracts both transverse directions. That contraction draws the rotating fluid inward toward the axis, which speeds up its interior rotation. Viscosity acts throughout this motion, so the spin-up gives a sharpening contribution without deciding the net change.

Call the unit axis direction \(e\), and let \(L(t)\) be the parcel length. If the strain is locally uniform across the parcel and aligned with this axis, write
\[
  S:=\frac12\bigl(\nabla u+(\nabla u)^{\mathsf T}\bigr),
  \qquad
  Se=\sigma(t)e,
  \qquad
  \sigma(t)>0,
  \qquad
  \frac{\dot L(t)}{L(t)}
  =
  e\cdot Se
  =
  \sigma(t).
\]

The axial rate removes transverse area at the same logarithmic rate. Incompressibility keeps the material volume \(\mathcal V\) fixed, so define \(A(t):=\mathcal V/L(t)\) and the area-equivalent radius \(r_{\mathrm{eff}}(t):=\sqrt{A(t)/\pi}\). Then
\[
  \nabla\cdot u=0,
  \qquad
  \frac{d}{dt}(A L)=0,
  \qquad
  \frac{\dot A}{A}=-\sigma(t),
  \qquad
  \frac{\dot r_{\mathrm{eff}}}{r_{\mathrm{eff}}}
  =
  -\frac{\sigma(t)}{2}.
\]

The transverse length scale closes at half the axial logarithmic rate, but parcel geometry has not decided whether vorticity grows while viscosity is acting. With \(\omega:=\nabla\times u\), evaluate the vorticity on the tracked material trajectory:
\[
  \frac{D\omega}{Dt}
  =
  S\omega+\nu\Delta\omega,
  \qquad
  \frac{D}{Dt}
  :=
  \partial_t+u\cdot\nabla,
\]
and under the alignment \(\omega=|\omega|e\),
\[
  S\omega
  =
  \sigma(t)\omega,
  \qquad
  \frac12\frac{D|\omega|^2}{Dt}
  =
  \sigma(t)|\omega|^2
  +
  \nu\,\omega\cdot\Delta\omega.
\]

The vorticity magnitude on that trajectory strengthens only where this complete pointwise right-hand side is positive; the viscous contribution itself has no fixed sign. Finite kinetic energy does not close that local question. The energy identity proved in \Cref{prop:ClassicalKineticEnergyIdentity} and the antisymmetric part of \(\nabla u\) give
\[
  \norm{u(t)}_2
  \leq
  \norm{u_0}_2
  <\infty,
  \qquad
  \left|\frac12\bigl(\nabla u-(\nabla u)^{\mathsf T}\bigr)\right|_{\mathrm F}
  =
  \frac{|\omega|}{\sqrt2}
  \leq
  |\nabla u|_{\mathrm F}.
\]
The energy bound leaves the local velocity gradient uncontrolled. This is enough to motivate a separate exact test of whether viscosity gains relative strength when the whole equation is viewed at finer scale.

\divider{\NavierStokes{} Scaling}

For \(\lambda>1\), form the rescaled field
\[
  u_\lambda(x,t)
  :=
  \lambda u(\lambda x,\lambda^2t),
  \qquad
  p_\lambda(x,t)
  :=
  \lambda^2p(\lambda x,\lambda^2t).
\]
This is another solution at another scale, not a later state of the tracked parcel. Each momentum term transforms with it:
\[
\begin{aligned}
  \partial_tu_\lambda(x,t)
  &=
  \lambda^3(\partial_tu)(\lambda x,\lambda^2t),
  \\
  (u_\lambda\cdot\nabla)u_\lambda(x,t)
  &=
  \lambda^3\bigl((u\cdot\nabla)u\bigr)(\lambda x,\lambda^2t),
  \\
  \nabla p_\lambda(x,t)
  &=
  \lambda^3(\nabla p)(\lambda x,\lambda^2t),
  \\
  \nu\Delta u_\lambda(x,t)
  &=
  \lambda^3\nu(\Delta u)(\lambda x,\lambda^2t).
\end{aligned}
\]
All four terms acquire the same cubic factor. Finer scale gives viscosity no automatic advantage over nonlinear transport, so the unresolved contest has to be located among quantities that scale by different powers.

\divider{Critical Height}

Let \(\Lambda:=(-\Delta)^{1/2}\). Since
\[
  \widehat{u_\lambda}(\xi,t)
  =
  \lambda^{-2}\widehat u(\lambda^{-1}\xi,\lambda^2t),
\]
a change of variables gives
\[
  \norm{u_\lambda(t)}_{\dot H^s}^2
  =
  \lambda^{2s-1}
  \norm{u(\lambda^2t)}_{\dot H^s}^2.
\]
At \(s=0\), kinetic energy loses one power of \(\lambda\); at \(s=1\), the first derivatives gain one. The exponent selects the unchanged level between them at \(s=\tfrac12\). Set
\[
  H(t)
  :=
  \frac12\norm{u(t)}_{\dot H^{1/2}}^2
  =
  \frac12\norm{\Lambda^{1/2}u(t)}_2^2.
\]
Writing \(H_\lambda\) for this global functional evaluated on \(u_\lambda\),
\[
  \norm{u_\lambda(t)}_2^2
  =
  \lambda^{-1}\norm{u(\lambda^2t)}_2^2,
  \qquad
  H_\lambda(t)=H(\lambda^2t),
  \qquad
  \norm{\nabla u_\lambda(t)}_2^2
  =
  \lambda\norm{\nabla u(\lambda^2t)}_2^2.
\]
Energy falls and the first-derivative norm rises under rescaling, while the global half-derivative height remains neutral. Scale neutrality says nothing about whether that height rises or falls along the original solution.

\divider{Nonlinear Production versus Viscous Dissipation}

Its nonlinear contribution is the signed production
\[
  \mathcal P_{1/2}[u](t)
  :=
  -\operatorname{Re}
  \pair
    {\Lambda^{1/2}u(t)}
    {\Lambda^{1/2}\bigl((u(t)\cdot\nabla)u(t)\bigr)}_{L^2}.
\]
The pressure pairing vanishes because \(\nabla\cdot u=0\), and the Laplacian contributes \(-\nu\norm{\Lambda^{3/2}u}_2^2\). Thus
\[
  H'(t)
  +
  \nu\norm{\Lambda^{3/2}u(t)}_2^2
  =
  \mathcal P_{1/2}[u](t).
\]
The height rises only when the signed nonlinear production exceeds the simultaneous positive viscous dissipation. Under the same rescaling,
\[
  \mathcal P_{1/2}[u_\lambda](t)
  =
  \lambda^2\mathcal P_{1/2}[u](\lambda^2t),
  \qquad
  \nu\norm{\Lambda^{3/2}u_\lambda(t)}_2^2
  =
  \lambda^2\nu\norm{\Lambda^{3/2}u(\lambda^2t)}_2^2.
\]
Both signed production and dissipation acquire the same square factor. Their critical balance still supplies no evolution interval for the original flow, leaving the viscous duration of a localized transverse scale to be found from a linear comparison.

\divider{The Heat Clock}

For the heat equation \(v_t=\nu\Delta v\),
\[
  \widehat v(\xi,t+\delta t)
  =
  e^{-\nu|\xi|^2\delta t}\widehat v(\xi,t).
\]
Now consider localized rotation or vorticity across a transverse width \(r\). Width alone does not locate its active frequencies. Suppose separately that a fixed positive fraction of its spectral mass lies in a band \(c r^{-1}\leq|\xi|\leq C r^{-1}\), with \(c\), \(C\), and that fraction independent of \(r\). On that band, viscous change becomes order one when
\[
  \nu|\xi|^2\delta t\simeq1,
\]
which gives the linear comparison time
\[
  \tau_\nu(r)
  :=
  \frac{r^2}{\nu}.
\]
This conditional clock is not a parcel lifetime or a nonlinear transition interval. It does show that halving the localized transverse width quarters the comparison time. More generally,
\[
  \tau_\nu(\lambda^{-1}r)
  =
  \lambda^{-2}\tau_\nu(r).
\]
The shorter clock now repeats geometrically, and its total length can be tested before asking whether one flow can realize it.

\divider{The Zeno Cascade}

For
\[
  r_n:=2^{-n}r_0,
  \qquad
  \tau_n:=\frac{r_n^2}{\nu},
\]
the heat clocks satisfy
\[
  \tau_n
  =
  \frac{r_0^2}{\nu}4^{-n},
  \qquad
  \sum_{n=0}^{\infty}\tau_n
  =
  \frac{4r_0^2}{3\nu}
  <
  \infty.
\]
The finite sum contains no material identity and no nonlinear evolution. If one \NavierStokes{} flow realizes this history, the same parcel passes through
\[
  r_0>r_1>r_2>\cdots\longrightarrow0
\]
at times
\[
  t_0<t_1<t_2<\cdots\uparrow T_*<\infty,
  \qquad
  t_{n+1}-t_n\asymp\tau_n,
\]
with comparison constants independent of \(n\). Its axial strain then accumulates along the actual trajectory so that \(\int_{t_0}^{t_n}\sigma(s)\,ds=2n\log 2+\varepsilon_n\), where \(|\varepsilon_n|\) is bounded independently of \(n\), and each actual next interval pays \(\int_{t_n}^{t_{n+1}}\sigma(s)\,ds=2\log 2+\delta_n\) with a uniform bound on \(|\delta_n|\).

At each \(t_n\), let \(\ell_n\) be the transverse localization width of the rotation. Positive constants \(c_0,C_0\), independent of \(n\), impose the mutual comparison \(c_0 r_n\leq r_{\mathrm{eff}}(t_n),\ell_n\leq C_0r_n\). The flow must also supply a localization or scale-transfer mechanism that carries one coherently traceable rotation through these widths; parcel contraction by itself does not perform that job.

Viscosity crosses the parcel boundary during every interval, so successive nearby snapshots need not descend from one localized rotation. Controlled diffusive exchange must preserve a coherent spacetime ancestry. Even that ancestry can persist while the amplitude fades. Choose cutoffs \(\chi_n\) adapted to the widths \(\ell_n\) and name the scale-normalized localized vorticity amplitude \(\mathfrak A_n:=r_n^{1/2}\|\chi_n\omega(t_n)\|_{L^2}\). The conditional history imposes \(\mathfrak A_n\geq a_0>0\), with \(a_0\) independent of \(n\); the signed vorticity balance is what this amplitude hypothesis has to survive, not a proof of the bound.

The shrinking localization widths still do not select an active band. A separate spectral-localization hypothesis places at least a fixed fraction \(\eta>0\) of \(\chi_n\omega(t_n)\) in \(c_1r_n^{-1}\leq|\xi|\leq C_1r_n^{-1}\), with \(c_1,C_1,\eta\) independent of \(n\). That active band makes the heat clock relevant, while the imposed gaps \(t_{n+1}-t_n\asymp r_n^2/\nu\) give the same flow the actual time in which to accumulate the next dyadic contraction. With these uniform gaps, the remaining time satisfies
\[
  T_*-t_n
  \asymp
  \sum_{k=n}^{\infty}\tau_k
  =
  \frac{4r_n^2}{3\nu}.
\]
Inside this conditional history, the same flow produces each next contraction in less time while diffusion continues: both the next interval and the total time left shrink like \(r_n^2/\nu\).
